%% ============================================================================
%%  pipeline.tex — THE HERO FIGURE
%%
%%  A fragment, not a document. \input it from poster.tex, or compile
%%  pipeline_preview.tex to iterate on it on its own (much faster).
%%
%%  Canvas is 114.3 cm wide = 45 in = the poster's \textwidth at 1.5 in margins,
%%  and ~29 cm tall. Coordinates are in cm at 1:1, so every font size below is a
%%  REAL point size on the printed poster. Do NOT wrap this in \resizebox — that
%%  would scale the type away from the sizes set in grinposter.sty.
%%
%%  Left to right: five numbered stages. The dashed backdrops split them into
%%  what happens ONCE, OFFLINE (stages 1-3) and what happens PER TRIAL, ONLINE
%%  (stages 4-5). That split IS the amortisation argument, so it is drawn first
%%  and given the most contrast.
%%
%%  LAYOUT BUDGET (x, cm):  band1 0.2-68.6 | gap | band2 70.6-114.2
%%    S1 frame  2.2-18.2      arrow 19.6-27.0
%%    S2 matrix 28.4-39.6     arrow 40.8-48.4
%%    S3 net    50.0-63.8     arrow 65.4-75.4  (crosses the offline/online line)
%%    S4 frame  77.0-93.0     arrow 94.4-99.0
%%    S5 tiles  98.8-113.3
%% ============================================================================

\usetikzlibrary{decorations.pathmorphing}

%% --- reusable pieces -------------------------------------------------------

% A GRT perceptual space: 4 stimulus distributions, axis-aligned (DS) bounds.
% Drawn inside a local 16 x 12 box with origin at bottom-left.
\newcommand{\grtframe}{%
  \draw[gmute, line width=1.6pt] (0,0) rectangle (16,12);
  % Decision bounds — axis aligned (that is DS). Their POSITION is free, and
  % that freedom is exactly what response bias is.
  \draw[greddeep, line width=2.2pt, dash pattern=on 7mm off 4mm] (8.4,0) -- (8.4,12);
  \draw[greddeep, line width=2.2pt, dash pattern=on 7mm off 4mm] (0,5.6) -- (16,5.6);
  \node[font=\fontsize{22}{22}\selectfont, text=gink] at (8,-0.95)
       {dimension A\ \ (e.g.\ border)};
  \node[font=\fontsize{22}{22}\selectfont, text=gink, rotate=90] at (-0.95,6)
       {dimension B\ \ (e.g.\ colour)};
}

% Single-contour version — the true, generating representation (stage 1).
\newcommand{\grtsingle}{%
  \foreach \cx/\cy/\rot in {4.4/2.9/0, 4.4/8.9/0, 12.2/2.9/34, 12.2/8.9/0}{%
    \draw[gbluedeep, line width=2.4pt, fill=gblue, fill opacity=0.45,
          rotate around={\rot:(\cx,\cy)}] (\cx,\cy) ellipse (2.15 and 1.55);
    \fill[greddeep] (\cx,\cy) circle (0.16);}
  \node[font=\fontsize{21}{21}\selectfont, text=greddeep] at (14.2,0.85) {$\rho \neq 0$};
}

% Nested-contour version — a posterior, with its uncertainty (stage 4).
\newcommand{\grtposterior}{%
  \foreach \cx/\cy/\rot in {4.4/2.9/0, 4.4/8.9/0, 12.2/2.9/34, 12.2/8.9/0}{%
    \foreach \s/\o in {1.34/0.13, 1.00/0.22, 0.66/0.38}{%
      \draw[gbluedeep, line width=1.4pt, fill=gblue, fill opacity=\o,
            rotate around={\rot:(\cx,\cy)}]
            (\cx,\cy) ellipse ({2.15*\s} and {1.55*\s});}
    \fill[greddeep] (\cx,\cy) circle (0.16);}
}

% A 4 x 4 identification confusion matrix. #1 = cell size in cm.
% Rows listed bottom-up so the diagonal runs the right way.
\newcommand{\cmgrid}[1]{%
  \foreach \row [count=\ri from 0] in {%
      {0.04,0.13,0.15,0.68},{0.10,0.05,0.68,0.17},%
      {0.14,0.66,0.05,0.15},{0.72,0.16,0.09,0.03}}{%
    \foreach \v [count=\ci from 0] in \row {%
      \fill[gbluedeep, opacity=\v] (\ci*#1, \ri*#1) rectangle ++(#1,#1);
      \draw[gmute, line width=0.7pt] (\ci*#1, \ri*#1) rectangle ++(#1,#1);}}
  \draw[gink, line width=1.6pt] (0,0) rectangle (4*#1,4*#1);
}

% A small multilayer perceptron.
\newcommand{\mlplayer}[4]{% #1 name  #2 x  #3 how many nodes  #4 y centre
  \foreach \i in {1,...,#3}{%
    \coordinate (#1-\i) at (#2, {#4 + (\i - (#3+1)/2)*1.6});}}

\newcommand{\mlpedges}[4]{% #1 name A  #2 count A  #3 name B  #4 count B
  \foreach \i in {1,...,#2}{\foreach \j in {1,...,#4}{%
    \draw[gmute, line width=0.5pt, opacity=0.5] (#1-\i) -- (#3-\j);}}}

\newcommand{\mlpnodes}[3]{% #1 name  #2 count  #3 colour
  \foreach \i in {1,...,#2}{%
    \filldraw[fill=#3, draw=gbluedeep, line width=1.2pt] (#1-\i) circle (0.5);}}

%% --- the figure ------------------------------------------------------------
\begin{tikzpicture}[x=1cm, y=1cm,
    flow/.style   ={-{Stealth[length=9mm,width=6.5mm]}, line width=3.2pt, draw=gbluedeep},
    loop/.style   ={-{Stealth[length=9mm,width=6.5mm]}, line width=3.8pt, draw=gred},
    flowlab/.style={font=\fontsize{21}{25}\selectfont, text=gink, align=center},
    stagett/.style={font=\fontsize{33}{38}\selectfont\bfseries, text=gink, anchor=west},
    stagecap/.style={font=\fontsize{24}{31}\selectfont, text=gink, align=left,
                     anchor=north west, text width=19.0cm},
    band/.style   ={draw=gmute, line width=1.6pt, dash pattern=on 9mm off 6mm,
                     rounded corners=14pt},
    bandlab/.style={font=\fontsize{27}{31}\selectfont\bfseries, text=gmute}]

  %% ---- the offline / online split, drawn first so everything sits on top ---
  \draw[band] (0.2,1.2) rectangle (68.6,24.4);
  \draw[band, draw=gred, opacity=0.85] (70.6,1.2) rectangle (114.2,24.4);
  \node[bandlab, anchor=west]            at (0.8,25.0)
        {OFFLINE\ \ \textbullet\ \ done ONCE\ \ \textbullet\ \ hours on one GPU};
  \node[bandlab, anchor=west, text=gred] at (71.2,25.0)
        {ONLINE\ \ \textbullet\ \ EVERY TRIAL\ \ \textbullet\ \ microseconds};

  %% ======================= STAGE 1 — the representation ===================
  \begin{scope}[shift={(2.2,11.4)}]
    \grtframe \grtsingle
  \end{scope}
  \node[circle,fill=gred,text=gpaper,minimum size=1.5cm,inner sep=0pt,
        font=\fontsize{34}{34}\selectfont\bfseries] at (2.9,9.2) {1};
  \node[stagett] at (4.0,9.2) {Sample a representation};
  \node[stagecap] at (2.2,7.8) {%
    Locations and noise of the four stimuli in perceptual space and where
    the decision bounds fall.};

  \draw[flow] (19.6,17.4) -- (27.0,17.4);
  \node[flowlab] at (23.3,19.5) {exact\\forward model};

  %% ======================= STAGE 2 — simulate =============================
  \begin{scope}[shift={(28.4,12.2)}]
    \cmgrid{2.8}
    \node[font=\fontsize{21}{21}\selectfont, text=gink] at (5.6,-0.95) {response};
    \node[font=\fontsize{21}{21}\selectfont, text=gink, rotate=90] at (-0.95,5.6) {stimulus};
    \node[font=\fontsize{22}{26}\selectfont, text=gred, align=center] at (5.6,11.6)
         {+ RT quantiles \emph{(optional)}};
  \end{scope}
  \node[circle,fill=gred,text=gpaper,minimum size=1.5cm,inner sep=0pt,
        font=\fontsize{34}{34}\selectfont\bfseries] at (29.1,9.2) {2};
  \node[stagett] at (30.2,9.2) {Simulate the data};
  \node[stagecap] at (28.4,7.8) {%
    Pushing $\theta$ through the exact forward model to produce matrices is cheap to generate \textbf{millions} of $(\theta,x)$ pairs.};

  \draw[flow] (40.8,17.4) -- (48.4,17.4);
  \node[flowlab] at (44.6,19.5) {$1.2\times10^{6}$\\pairs};

  %% ======================= STAGE 3 — train ================================
  \begin{scope}[shift={(50.0,17.4)}]
    \mlplayer{LA}{0}{5}{0}   \mlplayer{LB}{4.6}{7}{0}
    \mlplayer{LC}{9.2}{7}{0} \mlplayer{LD}{13.8}{4}{0}
    \mlpedges{LA}{5}{LB}{7} \mlpedges{LB}{7}{LC}{7} \mlpedges{LC}{7}{LD}{4}
    \mlpnodes{LA}{5}{gpaper} \mlpnodes{LB}{7}{gblue}
    \mlpnodes{LC}{7}{gblue}  \mlpnodes{LD}{4}{gred}
    \node[font=\fontsize{21}{21}\selectfont, text=gink] at (0,-6.2)    {data};
    \node[font=\fontsize{21}{21}\selectfont, text=gink] at (13.8,-6.2) {$q(\theta \mid x)$};
  \end{scope}
  \node[circle,fill=gred,text=gpaper,minimum size=1.5cm,inner sep=0pt,
        font=\fontsize{34}{34}\selectfont\bfseries] at (50.7,9.2) {3};
  \node[stagett] at (51.8,9.2) {Learn the inverse};
  \node[stagecap] at (50.0,7.8) {%
    Train a network to map data back to $\theta$ with a
    \textbf{full posterior}. Expensive, but paid \textbf{once}.};

  \draw[flow, draw=gred] (65.4,17.4) -- (75.4,17.4);
  \node[flowlab, text=gred, align=center] at (70.4,20.3)
       {\textbf{amortised:} trained once,\\reused for every participant};

  %% ======================= STAGE 4 — infer ================================
  \begin{scope}[shift={(77.0,11.4)}]
    \grtframe \grtposterior
  \end{scope}
  \node[circle,fill=gred,text=gpaper,minimum size=1.5cm,inner sep=0pt,
        font=\fontsize{34}{34}\selectfont\bfseries] at (77.7,9.2) {4};
  \node[stagett] at (78.8,9.2) {Infer, in microseconds};
  \node[stagecap, text width=16.0cm] at (77.0,7.8) {%
    A real matrix in, a calibrated posterior out in \chk{$\sim$3\,\textmu s}.
    MLE needs \chk{$\sim$90\,ms} and sometimes fails; this \textbf{always}
    answers.};

  \draw[loop] (93.8,17.4) -- (96.2,17.4);
  \node[flowlab, text=gred] at (95.0,19.8) {where is\\this person\\confused?};

  %% ======================= STAGE 5 — act ==================================
  \begin{scope}[shift={(97.6,12.4)}]
    % four stimuli of a 2 x 2 design; one is selected for the next trial
    \foreach \tx/\ty/\wob/\two in {0/5.7/0/0, 5.7/5.7/1/0, 0/0/0/1, 5.7/0/1/1}{
      \draw[gmute, line width=1.2pt, rounded corners=3pt]
            (\tx,\ty) rectangle ++(4.8,4.8);
      \ifnum\wob=1
        \draw[greddeep, line width=1.8pt, fill=gred, fill opacity=0.55,
              decorate, decoration={random steps, segment length=2.4mm, amplitude=0.9mm}]
              (\tx+2.4,\ty+2.4) circle (1.5);
      \else
        \draw[greddeep, line width=1.8pt, fill=gred, fill opacity=0.55]
              (\tx+2.4,\ty+2.4) circle (1.5);
      \fi
      \ifnum\two=1
        \fill[greddeep, opacity=0.75] (\tx+2.4,\ty+2.4) circle (0.78);
      \fi}
    % the one chosen for the next trial
    \draw[gbluedeep, line width=3.4pt, rounded corners=3pt] (5.7,0) rectangle ++(4.8,4.8);
    \node[font=\fontsize{21}{21}\selectfont\bfseries, text=gbluedeep] at (8.1,-0.95)
         {next trial};
    % the learner
    \begin{scope}[shift={(13.3,3.1)}]
      \filldraw[fill=gblue, draw=gbluedeep, line width=1.8pt] (0,1.5) circle (1.1);
      \filldraw[fill=gblue, draw=gbluedeep, line width=1.8pt]
        (-1.8,-1.8) .. controls (-1.8,0.3) and (1.8,0.3) .. (1.8,-1.8) -- cycle;
    \end{scope}
    \draw[flow, draw=gred, line width=2.6pt] (10.8,3.1) -- (11.4,3.1);
  \end{scope}
  \node[circle,fill=gred,text=gpaper,minimum size=1.5cm,inner sep=0pt,
        font=\fontsize{34}{34}\selectfont\bfseries] at (98.3,9.2) {5};
  \node[stagett] at (99.4,9.2) {Choose the next stimulus};
  \node[stagecap, text width=15.0cm] at (97.6,7.8) {%
    Present the stimulus that targets the confusion we just measured --- then
    re-infer, and go again.};

  %% ---- the closed loop ---------------------------------------------------
  \draw[loop, rounded corners=22pt]
        (104.8,4.6) -- (104.8,2.6) -- (85.0,2.6) -- (85.0,4.6);
  \node[font=\fontsize{25}{29}\selectfont\bfseries, text=gred, align=center] at (94.9,1.7)
       {every response updates the matrix, and the posterior moves with it};

\end{tikzpicture}
